\section*{Chapter \ref{ch:g6:shapes} --- Triangles and Quadrilaterals}

\begin{solution}{exo:g6:shapes:1}
Follow the method: \omterm{def:g6:lines:objects}{segment} $[AB]$ of $6$ cm, arc of \omterm{def:g3:shapes:shapes}{radius} $5$ cm
centered at $A$, arc of \omterm{def:g3:shapes:shapes}{radius} $4$ cm centered at $B$; the crossing
point is $C$.
\end{solution}

\begin{solution}{exo:g6:shapes:2}
Base $[EF]$ of $3$ cm, then two arcs of \omterm{def:g3:shapes:shapes}{radius} $5$ cm centered at $E$
and $F$; their crossing is $D$. The two base angles $\widehat{DEF}$
and $\widehat{DFE}$ are equal (about $72.5^\circ$ each) --- symmetry of
the \omterm{def:g6:shapes:triangles}{isosceles triangle} (\cref{ex:g6:symmetry:proves}).
\end{solution}

\begin{solution}{exo:g6:shapes:3}
\omterm{def:g6:lines:objects}{Segment} of $4$ cm, then two arcs of \omterm{def:g3:shapes:shapes}{radius} $4$ cm centered at its
ends. Every angle measures $60^\circ$.
\end{solution}

\begin{solution}{exo:g6:shapes:4}
(a) \omterm{def:g6:shapes:triangles}{equilateral} (hence also \omterm{def:g6:shapes:triangles}{isosceles}). \quad
(b) \omterm{def:g6:shapes:triangles}{isosceles} ($5 = 5$). \quad
(c) \omterm{def:g6:shapes:triangles}{right-angled} and \omterm{def:g6:shapes:triangles}{isosceles} (two equal angles of $45^\circ$ \omterm{def:g5:solids:def}{face} two
equal sides). \quad
(d) \omterm{def:g6:shapes:triangles}{equilateral} (three equal angles).
\end{solution}

\begin{solution}{exo:g6:shapes:5}
The hypotenuse is the side opposite the \omterm{def:g3:shapes:rightangle}{right angle}, always the longest
side. In $MNP$ \omterm{def:g6:shapes:triangles}{right-angled} at $N$, the hypotenuse is $[MP]$: it never
touches the \omterm{def:g5:solids:def}{vertex} of the \omterm{def:g3:shapes:rightangle}{right angle}, $N$.
\end{solution}

\begin{solution}{exo:g6:shapes:6}
\emph{1. True}: a square has four \omterm{def:g3:shapes:rightangle}{right angles}, which is the
definition of a rectangle.

\emph{2. False}: a $5 \times 3$ rectangle has unequal sides.

\emph{3. True}: a square has four equal sides, the definition of a
\omterm{def:g6:shapes:quadrilaterals}{rhombus}.

\emph{4. True}: see \cref{ex:g6:shapes:recognize}.
\end{solution}

\begin{solution}{exo:g6:shapes:7}
Construction with \omterm{def:g6:lines:perp}{perpendiculars}, as in \cref{ex:g6:shapes:program}.
Both diagonals measure about $5.8$ cm: the diagonals of a rectangle
have the same length (\cref{prop:g6:shapes:diagonals}).
\end{solution}

\begin{solution}{exo:g6:shapes:8}
Draw a \omterm{def:g6:lines:objects}{segment} of $6$ cm; its midpoint $O$; the \omterm{def:g6:lines:perp}{perpendicular} at $O$;
on it, two points at $2$ cm from $O$ on either side. Joining the four
endpoints gives the \omterm{def:g6:shapes:quadrilaterals}{rhombus} (all four \omterm{def:g5:solids:def}{vertices} are at distances
$3$ cm and $2$ cm from $O$ along \omterm{def:g6:lines:perp}{perpendicular lines}, so the four sides
are equal by symmetry).
\end{solution}

\begin{solution}{exo:g6:shapes:9}
One correct program:
1. draw a \omterm{def:g6:lines:objects}{segment} $[BC]$ with $BC = 4$ cm;
2. draw the arc of center $B$ and \omterm{def:g3:shapes:shapes}{radius} $6$ cm above $[BC]$;
3. draw the arc of center $C$ and \omterm{def:g3:shapes:shapes}{radius} $6$ cm above $[BC]$;
4. name $A$ the crossing point of the two arcs;
5. draw the \omterm{def:g6:lines:objects}{segments} $[AB]$ and $[AC]$.
\end{solution}

\begin{solution}{exo:g6:shapes:10}
Whatever the position of $M$ on the circle, the angle $\widehat{AMB}$
measures $90^\circ$: the triangle $AMB$ is \omterm{def:g6:shapes:triangles}{right-angled} at $M$
whenever $[AB]$ is a \omterm{def:g4:geometry:circle}{diameter}. The proof waits for
\cref{ch:g8:cosine}.
\end{solution}

\begin{solution}{exo:g6:shapes:11}
Diagonals that cut each other at their midpoint and have equal lengths:
that is the property list of the \emph{rectangle}
(\cref{prop:g6:shapes:diagonals}). Nothing forces the diagonals to be
\omterm{def:g6:lines:perp}{perpendicular}, so it need not be a \omterm{def:g6:shapes:quadrilaterals}{rhombus} or a square --- a rectangle
is the right answer (a square would satisfy it too, but is not
forced).
\end{solution}

\begin{solution}{exo:g6:shapes:12}
With $4 \times 4$ dots there are $3 \times 3$ grid cells.
Grid-aligned squares: sides $1$, $2$ or $3$ give $9 + 4 + 1 = 14$
squares. Tilted squares also exist: the small ones whose sides are
diagonals of $1 \times 1$ cells (like $(1,0)$, $(2,1)$, $(1,2)$,
$(0,1)$) --- four of them; and the larger ones whose sides are
diagonals of $2 \times 1$ rectangles, like $(1,0)$, $(3,1)$, $(2,3)$,
$(0,2)$ --- two of them. All four sides of such a square are equal
because they are diagonals of identical rectangles. Total:
$14 + 4 + 2 = 20$ squares. (Any count with a correct tilted example
and the $14$ aligned ones is a success at this level.)
\end{solution}

\begin{solution}{pb:g6:shapes:1}
\textbf{1.} The four sides all measure the same (about
$4.2$ cm) and the four angles are right: the shape is a
\emph{square}. (Equal diagonals give the \omterm{def:g3:shapes:rightangle}{right angles}, as in the
rectangle; \omterm{def:g6:lines:perp}{perpendicular} ones give the equal sides, as in the
\omterm{def:g6:shapes:quadrilaterals}{rhombus}; both at once give the square,
\cref{prop:g6:shapes:diagonals} read backwards.)

\textbf{2.} A \emph{\omterm{def:g6:shapes:quadrilaterals}{rhombus}}: four equal sides (about $4.7$ cm),
but no \omterm{def:g3:shapes:rightangle}{right angle} --- exactly the construction of
\cref{exo:g6:shapes:8}.

\textbf{3.} The four angles come out right but the sides come in
two different lengths: a \emph{rectangle} --- this is the
situation of \cref{exo:g6:shapes:11}.

\textbf{4.} No equal-length claim survives measuring and no
\omterm{def:g3:shapes:rightangle}{right angle} appears; but opposite sides are equal two by two and
look \omterm{def:g6:lines:perp}{parallel}. It is the general ``slanted rectangle'' whose
diagonals merely bisect each other --- the
\emph{parallelogram}, named and studied in \cref{ch:g7:central}.

\textbf{5.} Diagonals crossing at their common midpoint:
\begin{center}
\renewcommand{\arraystretch}{1.2}
\begin{tabular}{l|l}
diagonals\dots & shape \\
\hline
equal and \omterm{def:g6:lines:perp}{perpendicular} & square \\
\omterm{def:g6:lines:perp}{perpendicular} only & \omterm{def:g6:shapes:quadrilaterals}{rhombus} \\
equal only & rectangle \\
neither & the nameless one (see \cref{ch:g7:central}) \\
\end{tabular}
\end{center}

\textbf{6.} Each corner triangle has one leg along each
diagonal: \omterm{def:g3:division:half}{half} of $8$ cm and \omterm{def:g3:division:half}{half} of $5$ cm, so legs of $4$ cm
and $2.5$ cm, with a \omterm{def:g3:shapes:rightangle}{right angle} between them --- the same data
for all four. Triangles built from the same two legs and the
\omterm{def:g3:shapes:rightangle}{right angle} between them are exact copies of one another
(construct them: the compass settings are identical). Their four
hypotenuses --- the sides of the shape --- are therefore equal:
the shape is a \omterm{def:g6:shapes:quadrilaterals}{rhombus}.

\textbf{7.} Open the compass to the length of one side, then
walk it around the shape: if its point and pencil land exactly
on each pair of neighbouring corners, all four sides equal that
one opening. No ruler needed.

\textbf{8.} Program:
\begin{enumerate}
  \item Draw a \omterm{def:g6:lines:objects}{segment} $[AC]$ with $AC = 6$ cm and place its
        midpoint $O$ (measure $3$ cm from $A$).
  \item Draw the \omterm{def:g6:lines:perp}{perpendicular} to $(AC)$ at $O$.
  \item On it, place $B$ and $D$ at $3$ cm from $O$, one on each
        side.
  \item Join $A$, $B$, $C$, $D$.
\end{enumerate}
Diagonals equal ($6$ cm), \omterm{def:g6:lines:perp}{perpendicular}, bisecting each other:
by the dictionary, $ABCD$ is a square.

\textbf{9.} Measuring gives about $4.2$ cm: \emph{longer} than
$3$ cm --- the common guess ``\omterm{def:g3:division:half}{half} the diagonal'' is wrong. The
provable \omterm{def:g3:division:half}{half}: to go from corner $A$ to the next corner $B$, one
can walk the side $[AB]$ straight, or take the detour
$A \to O \to B$ through the center, costing $3 + 3 = 6$ cm; the
straight side is shorter, so the side measures less than $6$ cm.
That it is also \emph{more} than $3$ cm is what the ruler shows;
\cref{ch:g8:pythagoras} will prove it and give the exact value
($3$ cm multiplied by the square root of $2$, about $4.24$ cm).

\textbf{10.} The four corner triangles at $H$ now come in
\emph{two} matching pairs: two with legs $3$ cm and $2.5$ cm
(hypotenuses $[AB]$ and $[AD]$), two with legs $5$ cm and
$2.5$ cm (hypotenuses $[CB]$ and $[CD]$). So $AB = AD$ and
$CB = CD$: a \omterm{pb:g6:shapes:1}{kite} has two pairs of equal \emph{neighbouring}
sides (the \omterm{def:g6:shapes:quadrilaterals}{rhombus} is the special \omterm{pb:g6:shapes:1}{kite} where all four match).

\textbf{11.} \omterm{def:g4:geometry:polygon}{Quadrilateral}: $2$ diagonals. \omterm{def:g4:geometry:polygon}{Pentagon}: $5$.
\omterm{def:g4:geometry:polygon}{Hexagon}: $9$.

\textbf{12.} From one \omterm{def:g5:solids:def}{vertex} of the \omterm{def:g4:geometry:polygon}{hexagon}, diagonals go to
every \omterm{def:g5:solids:def}{vertex} except itself and its two neighbours: $6 - 3 = 3$
diagonals. Six \omterm{def:g5:solids:def}{vertices} launch $6 \times 3 = 18$ diagonal-ends
--- but each diagonal has been counted at \emph{both} of its
ends, once from each endpoint, exactly like a handshake counted
by both shakers (\cref{pb:g6:wholes:1}). True count:
$18 \div 2 = 9$. It matches the drawing.

\textbf{13.} Ten \omterm{def:g5:solids:def}{vertices}: $10 - 3 = 7$ from each, so
$10 \times 7 \div 2 = 35$ diagonals. Twenty \omterm{def:g5:solids:def}{vertices}:
$20 \times 17 \div 2 = 170$ diagonals.

\textbf{14.} Try candidates: $10$ \omterm{def:g5:solids:def}{vertices} give $35$ (question
13); $11$ \omterm{def:g5:solids:def}{vertices} give $11 \times 8 \div 2 = 44$. Found: the
\omterm{def:g4:geometry:polygon}{polygon} has $11$ \omterm{def:g5:solids:def}{vertices}.

\textbf{15.} A triangle has no diagonals at all: every \omterm{def:g5:solids:def}{vertex}'s
``non-neighbours'' are\,\dots\ nobody, since the other two
\omterm{def:g5:solids:def}{vertices} are both its neighbours. The method agrees:
$3 - 3 = 0$ from each \omterm{def:g5:solids:def}{vertex}, total $0$. Handshakes and
diagonals are the same count in different clothes: pairs chosen
from a group, each pair counted once --- mathematics recycles
its good ideas.
\end{solution}
